\documentclass[11pt]{article}
\usepackage[a4paper,margin=27mm]{geometry}
\usepackage{amsmath,amssymb,amsthm,mathtools}
\usepackage[T1]{fontenc}
\usepackage{lmodern,microtype,xurl}
\usepackage[colorlinks=true,linkcolor=blue,citecolor=blue,urlcolor=blue]{hyperref}
\newtheorem{theorem}{Theorem}
\newtheorem{proposition}{Proposition}
\newtheorem{lemma}{Lemma}
\theoremstyle{definition}
\newtheorem{assumption}{Assumption}
\newtheorem{remark}{Remark}
\newcommand{\F}{\mathbb F}
\newcommand{\Oh}{\mathcal O}
\title{Ordinary Polynomial Multiplication over Every Fixed Finite Field\\
  \large A strict logarithmic saving on fixed tapes}
\author{\'Arp\'ad Csizmadia\\
  \small\href{mailto:arpad.csizmadia.public@gmail.com}{\texttt{arpad.csizmadia.public@gmail.com}}}
\date{31 October 2026\\ \small Version 0.1.0-preprint. Apache-2.0.}
\hypersetup{pdfauthor={\'Arp\'ad Csizmadia},pdftitle={Ordinary Polynomial Multiplication over Every Fixed Finite Field}}
\begin{document}
\maketitle

\begin{abstract}
For each fixed finite field $\F_q$, ordinary polynomial multiplication
has deterministic bit complexity $O_q(n(\log(n+2))^{1-2^{-38}})$ on
a fixed finite-tape machine. For each fixed odd field $\F_{p^e}$, the
stronger bound $O_{p,e}(n(\log(n+2))^{9999/10000})$ follows from the
cyclotomic-host reduction proved here. The characteristic-two branch
and its fixed-extension corollary are supplied by the accompanying
controlled-tensor preprint~\cite{ControlledCRT}. The all-field bound
uses the larger of the two exponents and does not assert the sharper
odd-field exponent in characteristic two. Lean formalization remains
unfinished; full formalization is planned for a later update.
\end{abstract}
\paragraph{AI assistance and formal verification.}
This work was developed with substantial generative AI assistance in
mathematical exploration, algorithm and proof development, writing,
coding, and literature review. OpenAI Codex~\cite{CodexTool26}
assisted release preparation and source inspection. Historical
model versions were not consistently recorded; no complete
model-level attribution is claimed. Full Lean formalization is unfinished
and planned for a later update. Partial Lean work in the broader project
is not included here and does not verify these complete theorems.
Finite computational checks cover only their stated scopes. This
manuscript is a working preprint open for mathematical review.

\section{Prior work and source lineage}
For each fixed prime $p$, Harvey, van der Hoeven and Lecerf obtained
$O_p(n\log n\,8^{\log^*n})$ bit operations~\cite{HVHL17}.
Harvey and van der Hoeven improved this to the unconditional
$O_p(n\log n\,4^{\log^*n})$ benchmark on deterministic fixed tapes
and supplied the cyclotomic hosts used here~\cite{HVH19Cyclo}.
Fixed-degree basis reduction extends that comparison to every
individually fixed finite field. Their later uniform finite-field
$O(n\log q\log(n\log q))$ theorem assumes a Linnik constant
$L<1+2^{-1162}$~\cite[Theorem 1.2]{HVH22Finite}.
Harvey's October 2026 discussion identifies the fixed-field problem
as a separate issue from the new integer-multiplication claims
\cite[Section 4]{Harvey26Context}; it is context, not review of this work.

These developments established the algebraic and bit-complexity
framework on which this paper builds. The $n\log n$ scale is a
natural next benchmark for the unconditional fixed-field problem:
the earlier uniform theorem reaches that scale with its stated
number-theoretic assumption, while the unconditional cyclotomic
algorithm retains the iterated-logarithmic factor. Our claims fix
the field and address that remaining gap through the inherited hosts,
new transfer accounting and accompanying binary constructions.

The source lineage is explicit. OpenAI's finite-network and tape
interchange method~\cite{OpenAI26} is an antecedent. The selected odd
transfer imports the precise stage-two machine theorem
\cite{Stage2Machine26}, not an end-to-end integer multiplier.
The native supplement retains community chart and contributor
attributions~\cite{Community26}; related mixed-recursion and
exact-complex RAM work is credited separately~\cite{Sussman26}.
General fixed-matrix tensor acceleration in Alman's lookup-table
model~\cite{Alman23} is prior art, with a different access model.
The supplement uses classical integer arithmetic
\cite{SS71,BrentZimmermann10} and classical linear/quadratic
Clifford descriptions~\cite{DehaeneDeMoor03}; its specific
community-derived coordinate constructions retain their own credit.
These classical and community developments provide essential
antecedents to the construction.

The contribution recorded here is the odd-field host transfer and
its combination with the accompanying binary construction, under the
stated field and machine quantifiers. The companion baseline note
\cite{GeneralBaseline} instead combines the accessible binary route
with the odd proposition below to obtain $O_q(n\log(n+2))$.
These are related results with shared foundations.

\section{Model and scope}
The input consists of the ordinary coefficient lists of two polynomials
of degree less than $n$ over a fixed field $\F_q$, in a fixed polynomial
basis. The output is the full ordinary coefficient list of their
product. Time counts bit operations, including encoding, setup,
movement, decoding, and output, on a deterministic machine with a fixed
finite alphabet and a fixed finite number of sequential tapes. Constants
may depend on the fixed field. No growing-field uniformity, practical
cryptographic speedup, or combined workspace bound is claimed.

The logarithm is base two; changing its base changes only constants.
All statements concern ordinary univariate multiplication. Group
convolution, transformed products, integer multiplication, and
multivariate products are different problems and are not used as output
substitutes.

\section{The odd-characteristic theorem}
\begin{proposition}[Odd-field ordinary product]
\label{prop:odd}
For each fixed odd prime $p$ and positive integer $e$, ordinary
degree-$<n$ polynomial multiplication over $\F_{p^e}$ is deterministic
and takes
\[
  O_{p,e}\!\left(n(\log(n+2))^{9999/10000}\right)
\]
bit operations on a fixed finite-tape machine.
\end{proposition}
\begin{proof}[Proof source and dependency closure]
The detailed source proof is bundled at
\path{proofs/odd_grouped_native_transfer.tex} and
\path{proofs/odd_cyclotomic_host_transfer.tex}; exact source hashes
and copy edits are recorded in \path{audit.json}. The complete current
native construction, with the cited proof labels, is preserved in
\path{proofs/native_odd_source_snapshot.tex.txt}. A dependency map and the
remaining external import are documented in
\path{proofs/PROOF_SOURCES.md}. The proposition is not inferred from
an audit label or from a finite experiment.

The host is the unconditional fixed-prime cyclotomic construction of
Harvey and van der Hoeven~\cite{HVH19Cyclo}: it supplies $2n<N<4n$, admissible host
primes, a divisor, and a principal root in $N^{o(1)}$ setup time. The
coefficient host is enlarged by appending $O(\log\log N)$ unused host
primes only after the public generator degrees are frozen. With
$L_0=\lceil\log_2N\rceil$, $a=\lceil\log_2L_0\rceil$, and
$R_0=2^{a^2}$, the enlargement yields
\[
 \log\alpha'=O(a^4),\qquad m'D'=\Theta(N),\qquad
 1-D'/\alpha'=O(a^2/L_0^3),\quad m'=N/\alpha',\quad
 D'=\varphi(\alpha').
\]
The root lift $Y_{\rm old}\mapsto U^{\alpha'/\alpha}$ transports
the old root identities; new input coefficients instead use
$Y_{\rm new}\mapsto U$. This distinction preserves the ordinary
coefficient interpretation. Both coefficient CRT components are
reconstructed, the affine top map is inverted, and $2n-1<N$ rules out
cyclic aliasing.

The odd-characteristic formal compiler works over every commutative
$\F_p$-algebra, including nonreduced algebras. It chooses radix three
for $p\ne3$ and radix five for $p=3$, proves the root differences are
units by polynomial Bezout identities, and explicitly normalizes the
inverse by $t^{-1}$. Its recurrence gives a uniform quasilinear-size
$O_p(v\log^{16}(2v))$ circuit of depth $O_p(\log^2(2v))$ after counted
copy trees. The arbitrary-dirty replay uses a subtractive second read
to cancel $FUz$; it does not assume clean scratch. Six fixed-kernel
convolutions implement each exact-width Toeplitz chunk, with source,
target, and gate ports all inside the parent volume. The resulting
native schedule has $O_p(\log^{35}(2R))$ layers.

For the Newton factorization the sign
$\beta_j=(-1)^ju^{j(j-1)/2}$ is retained. On the growing coefficient
algebra, split-component unit selectors express each required shear as
two unit shears. A fixed five-call word supplies the unit shear over a
fixed quadratic field; its diagonal conjugations are explicitly priced
as growing-ring operations. No fixed-kernel theorem is applied to the
growing algebra. CRT characters are normalized by both inverse roots
and $(m')^{-1}$, and the complement component is charged separately.
Classical packed arithmetic is used independently for small
coefficient-ring products, so there is no recursive call to the target
multiplier.

The native supplier is the retained odd-ring phase/matching/CRT
construction bundled in \path{proofs/native_odd_source_snapshot.tex.txt}
(including the native dock, phase, selector, controlled operation,
CRT, and stage-two phase results), with the native parameter proof
used at the broader shape derived in the host transfer. The named
non-elementary movement dependency is
\path{IntegerMultBounds/MathlibTM2/M8/Stage2Machine/Bound.lean},
declaration \texttt{IntegerMultBounds.MathlibTM2.M8.Stage2Machine.Bound.uniform\_bound}
at commit \texttt{fefd699f9389af961046997796261333539e5909}~\cite{Stage2Machine26}; its
precise source is linked in the bundled
\path{proofs/PINNED_STAGE2_BINARY_MOVEMENT_BRIDGE_20261030.txt}.
The bridge derives the arbitrary-width adapter and records the exact
source hash. This is a named prior theorem, not a new unproved hypothesis. The final
cost is
\[
 O_p\!\left(NL_0^{\gamma_0}(\log L_0)^C\right)
   +N\operatorname{polylog}L_0+N^{o(1)},\qquad
 \gamma_0=1-\frac{619}{5000000}<\frac{9999}{10000};
\]
the strict exponent gap absorbs all fixed secondary logarithms.
Fixed extensions use a constant number of prime-field products and
fixed basis maps. The complete itemized audit is
\path{proofs/ODD_CYCLOTOMIC_FULL_ADVERSARIAL_AUDIT_20261030.txt}.
Its finite C++ witness checks bounded identities and examples only;
the asymptotic conclusion follows from the derivation and cited
interfaces, not from those samples.
\end{proof}

\begin{lemma}[Fixed binary-basis extension reduction]
\label{lem:binary-extension}
If ordinary multiplication over $\F_2$ on degree-$<n$ inputs has bit
cost $T(n)$, then for each fixed $s\ge1$, ordinary multiplication
over $\F_{2^s}$ in a fixed polynomial basis has bit cost
$O_s(T(n)+n)$ on the same fixed-tape model.
\end{lemma}
\begin{proof}
Fix an $\F_2$-basis $\omega_1,\ldots,\omega_s$ of $\F_{2^s}$.
Write the input coefficient sequences as
$f=\sum_{i=1}^s\omega_i f_i$ and $g=\sum_{j=1}^s\omega_jg_j$,
where $f_i,g_j\in\F_2[X]$ have degree less than $n$. Compute the
$s^2$ ordinary binary products $f_i g_j$. For each basis coordinate
$k$, combine their coefficients using the fixed structure constants
$\omega_i\omega_j=\sum_k c_{ijk}\omega_k$. The fixed basis encoding,
the $O(s^3n)$ additions/scalar maps, and the final coefficient
decoding cost $O_s(n)$. Since $s$ is fixed, the $s^2$ binary calls
cost $O_s(T(n)+n)$, proving the claim.
\end{proof}

\begin{theorem}[All-fixed-field bound]
\label{thm:all-fields}
For every fixed finite field $\F_q$, ordinary degree-$<n$
multiplication in a fixed polynomial basis takes
\[
 O_q\!\left(n(\log(n+2))^{1-2^{-38}}\right)
\]
deterministic bit operations on a fixed finite-tape machine.
\end{theorem}
\begin{proof}
Write $q=p^e$. If $p=2$, the controlled-tensor theorem and its fixed
extension corollary in \cite{ControlledCRT} give the stated bound.
If $p$ is odd, Proposition~\ref{prop:odd} gives the stronger exponent
$9999/10000$. Since
$9999/10000<1-2^{-38}$, its logarithmic factor is bounded by the one
displayed here for $n\ge2$; finitely many smaller lengths are absorbed
in the field-dependent constant. Both branches use ordinary
coefficient representations and the fixed finite-tape bit-time model.
All field and basis constants are independent of $n$.
\end{proof}

\begin{theorem}[Unconditional fixed-field baseline]
\label{thm:baseline}
Every fixed finite field has an unconditional ordinary multiplication
algorithm in $O_q(n\log(n+2))$ bit operations on the stated machine.
\end{theorem}
\begin{proof}
This follows immediately from Theorem~\ref{thm:all-fields}, because
$(\log(n+2))^{1-2^{-38}}=O(\log(n+2))$. The independent Pascal-transform
route for characteristic two is also recorded in
\cite{SimpleBinary}. The odd fields follow from Proposition~\ref{prop:odd}.
\end{proof}
\section{Verification and limitations}
The bundled historical 26-item adversarial derivation audit records
a bounded whole-host witness (24 products and 1,920 assertions over
$\F_3$ and $\F_5$); its implementation is not included here.
The recorded witness checks finite map boundaries; it does
not establish the growing tape algorithms or asymptotic bound. The
characteristic-two claim is imported from the accompanying
controlled-tensor preprint, whose result is stated with its own proof
scope. The general bound combines that result with the odd proposition.

Neither the odd proposition nor the combined theorem has an end-to-end
Lean proof in this package. Full Lean formalization is unfinished and
planned for a later update. The package records no independent
human review, publication acceptance, novelty certification, or
practical speedup.

\begingroup
\small
\raggedright
\begin{thebibliography}{99}
\bibitem{HVH19Cyclo}
David Harvey and Joris van der Hoeven.
\newblock Faster polynomial multiplication over finite fields using
cyclotomic coefficient rings.
\newblock \emph{Journal of Complexity}, 54:Article 101404, 2019.
\newblock \href{https://doi.org/10.1016/j.jco.2019.03.004}{doi:10.1016/j.jco.2019.03.004}.
\newblock \href{https://www.texmacs.org/joris/cyclomult/cyclomult.pdf}{Author manuscript}.

\bibitem{Stage2Machine26}
\newblock \texttt{platypii/integer-mult-bounds-lean},
\emph{Stage-two machine bound for equal-width interchange}, declaration
\texttt{IntegerMultBounds.MathlibTM2.M8.Stage2Machine.Bound.uniform\_bound},
in \path{IntegerMultBounds/MathlibTM2/M8/Stage2Machine/Bound.lean},
commit \texttt{fefd699f9389af961046997796261333539e5909}.
\newblock Source SHA-256: \texttt{97bb4fb5a883e0c90c53173355cf310ac56e2d28860679a13b2854f91cd2ad2a}.
\newblock \href{https://github.com/platypii/integer-mult-bounds-lean/blob/fefd699f9389af961046997796261333539e5909/IntegerMultBounds/MathlibTM2/M8/Stage2Machine/Bound.lean}{Pinned source}.

\bibitem{SimpleBinary}
\'Arp\'ad Csizmadia.
\newblock \emph{Exact Binary Polynomial Multiplication from Fast Pascal Transforms}.
\newblock Accompanying preprint in this collection; source:
\href{../02-simple-binary/main.tex}{\texttt{papers/02-simple-binary/main.tex}}.

\bibitem{ControlledCRT}
\'Arp\'ad Csizmadia.
\newblock \emph{Ordinary Binary Polynomial Multiplication below $n\log n$}.
\newblock Subtitle: Controlled tensor routing on fixed tapes.
\newblock Accompanying preprint in this collection; source:
\href{../03-best-binary/main.tex}{\texttt{papers/03-best-binary/main.tex}}.
\bibitem{HVHL17}
David Harvey, Joris van der Hoeven, and Gr\'egoire Lecerf.
\newblock Faster polynomial multiplication over finite fields.
\newblock \emph{Journal of the ACM}, 63(6), Article 52, 2017.
\newblock \url{https://doi.org/10.1145/3005344}.
\newblock Author manuscript: \url{https://www.texmacs.org/joris/ffmul/ffmul.pdf}.

\bibitem{HVH22Finite}
David Harvey and Joris van der Hoeven.
\newblock Polynomial multiplication over finite fields in time $O(n\log n)$.
\newblock \emph{Journal of the ACM}, 69(2), Article 12, 1--40, 2022.
\newblock Theorem 1.2 assumes the stated Linnik hypothesis.
\newblock \url{https://doi.org/10.1145/3505584}.
\newblock Author manuscript: \url{https://www.texmacs.org/joris/ffnlogn/ffnlogn.pdf}.

\bibitem{Harvey26Context}
David Harvey.
\newblock \emph{On the recently proposed $o(n\log n)$ integer multiplication algorithm}.
\newblock Unpublished discussion note, 20 October 2026, Section 4.
\newblock \url{https://web.maths.unsw.edu.au/~davidharvey/research/subnlogn.pdf}.

\bibitem{OpenAI26}
OpenAI.
\newblock \emph{Integer multiplication below $n\log n$}.
\newblock Research preprint, 23 September 2026, Sections 3--4;
Lemma 12 supplies the arbitrary-width equal-chunk interchange statement.
\newblock Citation revision \texttt{4eaee87548250b2682affda7d05d947c05554bae}, inspected 31 October 2026.
\newblock \url{https://github.com/openai/math/blob/4eaee87548250b2682affda7d05d947c05554bae/preprints/Integer-multiplication-below-n-log-n-September-23-2026/build/sections/04-swap.tex}.

\bibitem{Community26}
Douglas Colkitt and the contributors to \texttt{CrocSwap/integer-mult-bounds}.
\newblock \emph{A sharper exponent for integer multiplication}.
\newblock Research repository, revision \texttt{2485805a6d47f36ea83ddde6327d453ac82d44dc},
accessed 31 October 2026. Contributor-level source notices are retained
in the bundled native derivation. Its selected integer bound retains
its stated framework assumptions.
\newblock \url{https://github.com/CrocSwap/integer-mult-bounds/tree/2485805a6d47f36ea83ddde6327d453ac82d44dc}.

\bibitem{Sussman26}
J.~A.~Sussman.
\newblock \emph{Walsh--Hadamard power saving: Lean proofs}.
\newblock Research repository, accessed 31 October 2026.
Related exact-complex RAM constructions and recursion ideas; its
end-to-end bounds are not premises of the selected fixed-field theorem.
\newblock \url{https://github.com/jacobalansussman/wht-power-saving-lean}.

\bibitem{Alman23}
Josh Alman.
\newblock Faster Walsh-Hadamard transform and matrix multiplication over finite fields using lookup tables.
\newblock In \emph{Symposium on Simplicity in Algorithms}, 137--144, 2023.
\newblock \url{https://doi.org/10.1137/1.9781611977585.ch13}.
\newblock \url{https://arxiv.org/abs/2211.04643}.

\bibitem{SS71}
Arnold Sch\"onhage and Volker Strassen.
\newblock Schnelle Multiplikation gro{\ss}er Zahlen.
\newblock \emph{Computing}, 7:281--292, 1971.
\newblock \url{https://doi.org/10.1007/BF02242355}.

\bibitem{BrentZimmermann10}
Richard P. Brent and Paul Zimmermann.
\newblock \emph{Modern Computer Arithmetic}.
\newblock Cambridge University Press, 2010.
Author-hosted version 0.5.9, Sections 1.3.5, 1.4.3 and 2.5.
\newblock \url{https://members.loria.fr/PZimmermann/mca/mca-cup-0.5.9.pdf}.

\bibitem{DehaeneDeMoor03}
Jeroen Dehaene and Bart De Moor.
\newblock The Clifford group, stabilizer states, and linear and quadratic operations over GF(2).
\newblock \emph{Physical Review A}, 68:042318, 2003.
\newblock \url{https://doi.org/10.1103/PhysRevA.68.042318}.

\bibitem{GeneralBaseline}
\'Arp\'ad Csizmadia.
\newblock \emph{Ordinary Polynomial Multiplication over Fixed Finite Fields in $O(n\log n)$ Time}.
\newblock Accompanying baseline note in this collection, 2026;
source and proof status in \path{papers/04-general-nlogn/}.

\bibitem{CodexTool26}
Nicolas Bonamy and Derrick Choi.
\newblock Codex as a platform: build on the open agent harness.
\newblock OpenAI Developers, 19 August 2026. Official documentation
reference for the Codex software used in release preparation and source
inspection; accessed 31 October 2026.
\newblock \url{https://developers.openai.com/blog/codex-as-a-platform}.

\end{thebibliography}
\endgroup

\end{document}
